% 제1장 서론
% 참조: doctoral/research_model/research_model.tex, doctoral/ontology/ontology_analysis.tex
% v9.3 기준값: qa/v9.3_band_seed/reference_values.json
% κ = 1.2, R₀ = 6.0h, R_max = 24.0h (= 4R₀)

\chapter{서론}
\label{ch:introduction}

\section{연구 배경}
\label{sec:background}

현대 사회는 자연재해, 사이버공격, 공급망 중단, 감염병 대유행 등 다양한 재난 및
업무 중단 위험에 직면하고 있다. 이러한 위험은 단일 사건으로 발생하기보다는
복합적으로 상호작용하며, 조직의 핵심 업무 프로세스에 심각한 영향을 미친다.
BCMS(Business Continuity Management System, 업무연속성관리체계)는 이러한 재난
상황에서 조직의 핵심 기능을 유지하고 신속하게 복구하기 위한 체계적 관리
프레임워크로서, \nrev{ISO 22301:2019 국제표준}을 중심으로
전 세계적으로 확산되고 있다.

BCMS의 수립 과정에서 BIA(Business Impact Analysis, 업무영향분석)은 핵심
구성 요소로 기능하며, 이를 통해 MTPD(Maximum Tolerable
Period of Disruption, 최대허용중단기간)과 RTO(Recovery Time Objective, 복구목표시간)이 결정된다
\citep{hiles2010bcm, snedaker2013bcdr}.
RTO는 재난 발생 후 업무 프로세스가 수용 가능한 수준으로 복구되어야 하는
목표 시간을 의미하며, $\text{RTO} \leq \text{MTPD}$의 관계가 성립한다
\nrev{(ISO 22313, 2020)}.
현행 \nrev{ISO 22301:2019 국제표준} 및 대부분의 BCM 실무에서 RTO는 BIA 단계에서 고정값으로
설정되며, BIA 주기 사이의 평상 운영 중에는 실시간으로 변경되지
않는다~\citep{herbane2010disaster}.

\chg{그러나 복구가 이루어지는 실제 운영 환경은 BIA 시점에 고정한
가정과 달리 끊임없이 변동하므로, 한 번 설정한 RTO는 실제 복구 상황과 어긋날
수 있다. 이러한 한계를 다루기 위해 \citet{lee2026drto}은 조직의
장애를 탐지하고 상황을 파악하는 능력인 관측성(observability)을 반영하여
RTO를 동적으로 조정하는 DRTO 모델을 제안하였다.

이 모델은 평상시 기준 복구시간 $R_0$에서 출발하여, 관측성의 효과를 시그모이드
함수로 변환해 기준 시간에서 차감하는 형태이다. 관측성이 높을수록 목표 시간을
더 많이 단축하되, 시그모이드 바운딩에 의해 복구시간이 $0$ 아래로 내려가지
않으며, 곡률은 매개변수 $\kappa$로 조절한다. 해당 연구에서는 정적 기준(C0)에서
관측성을 활성화한 조건(C1)으로의 전환을 설정하고, 대규모 몬테카를로
시뮬레이션과 비중첩 대역 시드 체계로 그 효과를 검증함으로써, 고정된 정적 RTO를
상황에 반응하는 동적 모델로 전환하는 진전을 이루었다.

그러나 이 모델은 관측성이라는 단일 차원에 머물렀다. 복구에 걸리는 시간은
장비 가용성과 인력 동원, 외부 공급망 등 여러 요인의 불확실성(uncertainty)
아래 놓인다. \chg{이 불확실성은 일상적 변동을 나타내는 가우시안(정규) 분포뿐
아니라, 그 범위를 크게 벗어나는 극단 사건을 나타내는 파레토(중꼬리) 분포까지를
포함한다.} 관측성에 국한된 모델은 이러한 불확실성을 반영하지 못한다. \chg{관측성이
복구시간을 단축하는 방향으로 작용한 것과 달리, 불확실성은 복구시간을 늘려
MTPD에 근접시키는 방향으로 작용하며, 특히 극단을 나타내는
파레토 영역에서 그 증가 폭이 커진다.} 본 연구는 이 관측성 모델에 불확실성
차원을 더하여 모델을 완성하는 데에서 출발한다. 다음 절에서는 관측성과 불확실성의
두 측면에서 현행 접근의 한계를 선행연구에 비추어 구체화하되, 불확실성은 다시
가우시안과 파레토의 두 종류로 나누어 살펴보고, 이를 통합적으로 해소할 필요성을
논증한다.}


\section{연구 필요성}
\label{sec:problem_statement}

\chg{연구 배경에서 제기한 관측성과 불확실성의 두 측면은 업무연속성 연구에서
개별적으로는 다루어져 왔으나, 그동안의 연구 및 실무에서는 RTO를 결정하는 하나의 틀 안에서
통합적으로 반영되지는 못하였다. 그 결과 정적 RTO는 두 측면 모두에서 현실과
어긋나며, 이는 복구 자원의 과부족과 목표 시간의 신뢰성 저하로 이어진다.
따라서 관측성의 한계를 먼저 검토한 뒤, 불확실성을 일상적 변동의 가우시안과
극단 사건의 파레토라는 두 종류로 나누어 선행연구에서 없었던 방법으로 구체화하고, 이들이
결합될 때 개별적 보정만으로는 해결되지 않는다는 점에서 RTO 산정에 대한
통합적 접근의 필요성을 다음과 같이 세 가지로 제시한다.}

\chg{첫째, 일상적 변동을 나타내는 가우시안과 극단 사건을 나타내는 파레토의
두 불확실성을 관측성 모델에 통합할 필요가 있다. 관측성은
\citet{lee2026drto}의 동적 모델 연구에서 이미 반영한 차원이다. 이 모델은 관측성이 높을수록 이상 징후를 조기에 탐지하여
복구를 가속할 수 있다는 점을 시그모이드 바운딩으로 포착함으로써, 정적 RTO가
야기하던 관측성 높은 조건의 과다 예비(over-provisioning)와 낮은 조건의 과소
예비(under-provisioning)를 완화하였다~\citep{majors2022observability}. 그러나
이 모델은 관측성이라는 단일 차원에 국한되어, 복구 과정의 불확실성을 통합하지
못하였다. 이 불확실성은 일상적 변동을 나타내는 가우시안과 극단성을 나타내는
파레토의 두 가지로 나타나므로, 이 둘을 하나의 모델로 통합하여 반영하는 것이
본 연구가 해소하고자 하는 첫째 필요성이 된다.

둘째, 불확실성 정량화 부재 문제이다. 복구 과정에는 장비 가용성과 인력 동원
속도, 외부 공급망 의존도, 규제 대응 절차 등 다양한 요인이 작용하며, \chg{이들이
복구 시간에 미치는 영향은 재난의 종류에 따라 크게 달라진다. 예컨대 사이버공격은
데이터 복원과 시스템 재구성에, 자연재해는 물리적 설비와 인력 접근성에, 감염병
대유행은 인력 가용성과 외부 의존성에 각기 다른 부담을 지운다. 이처럼 같은
조직이라도 재난 유형과 상황에 따라 복구 소요가 달라지므로, 이 요인들은
본질적으로 확률적으로 변동한다.} 그럼에도 기존 RTO 프레임워크는 이러한
불확실성을 ``최악의 경우(worst case)''나 ``평균적 시나리오''라는 하나의
시나리오로 단순화할 뿐, 확률 분포에 기반한 정량화를 수행하지
않는다~\citep{cerullo2004bcp, gibb2006bcm}. 그 결과 BIA에서 설정한 RTO와 실제
복구 시간 사이에 구조적 괴리가 생기고, 조직은 그 괴리가 어느 방향으로 얼마나
벌어질지를 사전에 가늠할 수 없다. 다수 요인의 합으로 근사되는 일상적 복구
변동은 가우시안(정규) 분포로 표현될 수 있으므로, 불확실성을 확률 분포로
명시하면 RTO를 단일 점 추정이 아니라 분포의 함수로 다룰 수 있다. 이 차원은
관측성 기반 모델에서는 다루어지지 않았으므로, 일상적 복구 변동의 불확실성을
확률 분포로 정량화하여 복구목표시간 모형에 반영하는 일이 요구된다.

셋째, 극단 사건 미모형화 문제이다. 이는 앞의 일상적 변동과 더불어 불확실성이
갖는 또 다른 형태로, 불확실성을 확률 분포로 다루더라도 대부분의 접근이
정규분포만을 전제하는 데에서 비롯된다~\citep{coles2001extremes, embrechts1997extremal, taleb2007blackswan}.
그러나 사이버공격에 의한 데이터 센터 전면 장애, 복합 재난에 따른 공급망 동시
중단, 핵심 인력의 대규모 이동 제한과 같은 극단 사건은 정규분포의 가정을 크게
벗어난다. 이러한 사건의 복구 시간은 평균에서 멀어질수록 발생 확률이 급격히
줄지 않는 중꼬리(heavy-tail) 특성을 보이며, 극단값이론(Extreme Value Theory, EVT)에 따르면 파레토
분포가 이를 적절히 기술한다~\citep{coles2001extremes, embrechts1997extremal}. 중꼬리 분포에서는 고분위 영역의 복구 시간이
정규분포에 비해 현저히 커지므로, 정규분포를 전제한 RTO 산정은 극단 상황에서
실제 소요 시간을 크게 과소평가한다. 정적 RTO와 관측성 단일 모델은 모두 이 꼬리
위험(tail risk)을 구조적으로 포착하지 못하여, 업무연속성이 가장 절실한 극단
상황에서 보호에 실패할 수 있다. 따라서 정규분포의 범위를 넘어, 극단값이론에
근거하여 극단 사건을 중꼬리 분포로 별도로 모형화하고 그 꼬리 위험을
RTO에 반영하는 일이 요구된다.}

\nrev{\chg{나아가 관측성과 불확실성은 서로 독립적인(직교하는) 두 차원으로서
복구 시간에 서로 반대 방향으로 작용한다. 관측성이 높을수록 이상 징후를 일찍
탐지하여 복구 시간을 단축하는 반면, 불확실성이 클수록 복구 시간의 변동과 지연을
키운다. 특히 불확실성의 분포가 정규분포를 벗어나 중꼬리(파레토) 특성을 가지면
극단 영역에서 복구 시간이 급증하므로, 정규분포를 전제한 통상적 위험 평가는 그
타당성을 잃는다. 두 차원은 서로 독립적이더라도 복구 시간이라는 하나의 결과
위에서 함께 작용하므로, 각 차원을 따로 다루어서는 하나의 동적복구목표시간에 온전히
반영하기 어렵다.

이때 각 요인을 서로 분리하여 독립적·선형 보정량으로 추정한 뒤 경계 없이 단순
합산하는 접근(decoupled approach)은 다음과 같은 구조적 한계를 지니며, 바로
여기에서 본 연구의 필요성이 비롯된다. 첫째, 보정량을 경계 없이 더하는 방식은
복구 시간이 음수가 되거나 최대허용중단기간을 넘어서는 등 물리적으로 타당한
범위를 벗어날 수 있다. 둘째, 실무에서는 동적복구목표시간을 하나의 시간값으로
결정해야 하므로, 여러 보정 계수를 사후에 합산하는 구조는 의사결정 절차상
비현실적이며 책임 소재가 분산된다. 셋째, 각 차원을 서로 다른 척도와 절차로
분리해 보정하면, 관측성의 단축 효과와 불확실성의 연장 효과를 하나의 복구 시간
위에서 일관되게 비교하고 결합하기 어렵다. 요컨대 독립적·선형·무경계 보정의 단순 합산만으로는 두 차원을 물리적
경계 안에서 하나의 동적복구목표시간으로 통합할 수 없으며, 여기에서 통합적 접근의
필요성이 도출된다.}

\chg{다음 절에서는 이러한 필요성으로부터 도출되는 본 연구의 목적을
구체화한다.}}


\section{연구 목적}
\label{sec:research_objectives}

\chg{본 연구는 \citet{lee2026drto}이 제안한 관측성 기반
동적 복구목표시간($\DRTO$) 모델을 불확실성까지 반영하는 모델로 확장하고 그
타당성을 검증하는 것을 목적으로 한다. 기존 모델은 기준 복구시간에 관측성의
효과를 더하는 단일 채널 구조로, 관측성을 시그모이드 바운딩으로 반영하여
복구시간을 물리적으로 타당한 범위 안에 유지한다. 본 연구는 동일한 바운딩
원리를 불확실성에 적용하여, 관측성과 불확실성의 이중 채널 모델로 격상한다.
확장된 모델의 구체적 수리 구조와 도출 근거의 이론적 배경은
\chref{ch:theory}에서 설명한다. 이러한 세 가지 필요성(한계)에
대응하는 구체적 연구 목적은 다음과 같이 다섯 가지로 설정하였다.

첫째, 관측성 단일 채널 모델을 관측성과 불확실성의 이중 채널 $\DRTO$ 모델로
확장한다. 관측성을 복구 능력의 척도로 반영하고(관측성 미반영 문제 대응),
불확실성을 가우시안 분포로 정량화하며(정량화 부재 문제 대응), 일상적 변동을
넘는 극단을 파레토 분포로 포착한다(극단 미모형화 문제 대응). 두 채널은 동일한
시그모이드 바운딩 원리를 공유하므로, 불확실성 채널을 더한 뒤에도 복구시간이
물리적으로 타당한 범위를 벗어나지 않는 일관된 구조를 유지한다. 이로써 세 한계를
하나의 모형 안에서 동시에 해소한다.

둘째, 정적 기준에서 관측성과 불확실성을 단계적으로 도입하는 네 개 운영 조건
체계(C0--C3)를 설계한다. 이 조건 체계는 i) 정적 기준(C0), ii) 관측성(C1)
기준, iii) 관측성(C1)과 가우시안 불확실성(C2)을 함께 적용하는 기준, 그리고
iv) 관측성(C1)과 파레토 불확실성(C3)을 함께 적용하는 기준을 차례로 하여 각
요인이 추가될 때의 한계 기여를 분리하여 관측함으로써, 요인별 효과와 결합 효과를
명확히 귀속하고 비교할 수 있는 분석 틀이 된다.

셋째, 확장된 모델의 타당성을 검증할 연구 가설을 수립하고 체계적으로 검증한다.
가설은 모델의 핵심 효과와 결과의 강건성, 회귀 설명력, 이중채널 비선형 효과,
전문가 타당성의 다섯 측면을 포괄하며, 이를 대규모 몬테카를로 시뮬레이션과
전문가 평가를 결합한 다층 검증으로 확인한다.

넷째, 가우시안 불확실성(C2)과 파레토 불확실성(C3)을 각각 독립적으로 적용한
결과를 분포 전 구간에 걸쳐 비교하여, 극단 영역에서 나타나는 이중채널 비선형
감쇠 메커니즘을 규명하고 정량화한다. 두 조건은 평균 수준에서는 거의 구분되지
않으나 특정 백분위(약 P85.8)를 경계로 분포가 역전되어, 극단 영역에서는 파레토
불확실성이 복구 시간을 급격히 연장하는 한편 관측성 채널의 한계 효과는 그
영역에서 감쇠한다. 이러한 핵심 영역과 꼬리 영역 사이의 비선형 거동은 두 조건을
분포 전 구간에서 함께 비교·분석할 때에만 드러나며, 이는 확장 모델이 갖는 고유한
설명력에 해당한다.

다섯째, 검증된 모델을 바탕으로 제조업, 금융, 의료, 정보기술 서비스, 물류,
소매 등 주요 산업 분야에 적용할 수 있는 실무 가이드라인을 제시한다. 산업마다
허용 가능한 복구시간과 위험 노출이 다르므로, 가이드라인은 산업별 특성에 맞추어
모델의 매개변수를 조정하는 방식으로 제공한다.

\chg{요컨대 본 연구는 관측성에 국한된 기존 동적 복구목표시간 모델을 관측성과
불확실성의 이중 채널 모델로 확장하고, 이를 검증할 네 개 운영 조건 체계와 연구
가설을 수립하여 다층 검증으로 그 타당성과 극단 영역의 이중채널 비선형 감쇠를
확인하며, 검증된 모델을 산업별 실무 가이드라인으로 구체화하는 것을 목적으로
한다.}

이상의 목적에서 확장 모델의 수리적 구조와 매개변수, 조건 체계, 연구 가설의
도출은 \chref{ch:theory} 이론적 배경에서, 검증 절차와 다층 검증 체계는
\chref{ch:method} 연구 방법에서 선행 이론에 근거하여 상세히 전개한다.}


\section{연구 범위 및 방법}
\label{sec:research_scope}

\chg{이러한 목적을 달성하기 위해 본 연구의 범위와 방법을 다음과 같이 정한다.
모델의 범위는 정적 기준(C0)에서 관측성(C1), 가우시안 불확실성(C2), 파레토
불확실성(C3)에 이르는 네 개 운영 조건을 대상으로 한다. 이 네 조건은 앞서
설정한 이중 채널 모델의 구성 요소가 하나씩 활성화되는 단계에 대응하며,
복잡성이 점증하는 순서로 배열된다. 확장된 $\DRTO$ 모델은 관측성과 불확실성을
개별 시그모이드 바운딩으로 결합하는 구조이며, 그 수리적 정의와 매개변수 설정
근거는 \chref{ch:theory} 이론적 배경에서 전개한다.

검증 방법으로는 \citet{lee2026drto}에서 정립한 몬테카를로 시뮬레이션과
비중첩 대역 시드(band seeding) 체계를 따른다. 모델이 확률 분포를 입력으로
받으므로, 그 결과인 $\DRTO$의 분포를 하나의 공식으로 직접 구하기는
어렵다. 따라서 난수 표본을 많이 생성하여 그 결과를 통계적으로 추정하는
몬테카를로 방법이 적합하다. 본 연구는 이 방법론을 확장한 조건 체계에
적용하여, $n=10{,}000$ 규모의 시뮬레이션으로 조건별로 모델이 산출하는 결과를
통계적으로 확인하고, 여기에 업무연속성관리 및 재해복구
분야 전문가 평가를 결합하여 그 실무 타당성을 보완적으로 검증한다. 특히
시뮬레이션과 회귀, 전문가 평가 등 서로 다른 방법의 결과가 동일한 결론으로
수렴하는지를 확인하는 삼각검증(triangulation)을 통해 결론의 타당성을 높인다.
\chg{이러한 검증은 여섯 층위로 이루어진 다층 검증 체계(L0--L5)로 구조화된다.
먼저 기반 검증(L0)에서 모델 구현의 정합성과 물리적 경계를 결정론적으로
확인하고, 시뮬레이션 검증(L1)에서 관측성 효과와 경계 보장 등 모델의 기본
작동을 확인한다. 이어 회귀 검증(L2)으로 변수 간 설명력과 교호작용을, 강건성
검증(L3)으로 결과의 재현성을, 이중채널 검증(L4)으로 관측성과 극단 불확실성의
비선형 상호작용을 살핀 뒤, 전문가 평가(L5)로 실무 타당성을 확인한다. 각 층위의
구체적 정의와 절차는 \chref{ch:method} 연구 방법에서 제시한다.} 이를 통해
모델의 핵심 효과와 강건성, 실무 타당성을 포괄하는 연구 가설을 검증하며, 그
가설의 도출 근거는 \chref{ch:theory} 이론적 배경에서 상세히 기술한다.

한편 본 연구는 이론적 모델 수립과 시뮬레이션 검증, 전문가 타당성 평가에
초점을 맞추며, 다음 세 가지는 범위에 포함하지 않는다. 즉 i) 실제 조직
환경에서의 현장 실험(field test), ii) 시간에 따라 RTO가 연속적으로 변화하는
시계열 DRTO 모델, iii) 실제 단일 조직에 대한 심층 사례 연구(in-depth
single-case study)이다. 이들은 모델의 이론적 타당성과 여러 상황에서의 일반적
작동을 먼저 확립한다는 본 연구의 목적에 비추어 후속 과제로 두며, 검증된 모델은 이러한
현장 적용 연구의 토대가 된다.}



\section{논문 구성}
\label{sec:thesis_structure}

본 논문은 총 5개 장과 부록으로 구성되며, 각 장의 내용은 다음과 같다.

\vrev{%
\textbf{\chref{ch:introduction} 서론}에서는 연구 배경, 연구 필요성, 연구 목적, 연구 범위 및
방법, 논문 구성을 서술한다.

\textbf{\chref{ch:theory} 이론적 배경}에서는 업무연속성관리체계(BCMS)와
복구목표시간(RTO), 시스템 관측성 이론, 불확실성 정량화, 극단값
이론(EVT), 시그모이드 바운딩 기법에 관한 기존 연구를 검토하고
연구 격차를 도출한다. 이어서 개별 시그모이드 바운딩 모델의
수리적 정의, 4단계 조건 체계(C0--C3), 매개변수 설정($R_0 = 6.0$h,
$R_{\max} = 24.0$h, $\kappa = 1.2$), 시뮬레이션 설계 요소(비중첩
대역 시드 체계), 그리고 \jrev{10개} 가설(H1--\jrev{H10})을 수립한다.

\textbf{\chref{ch:method} 연구 방법}에서는 4단계 조건 체계(C0--C3)에 대응하는
연구 설계와 $n = 10{,}000$ 몬테카를로 시뮬레이션의 실행 절차,
분포 매개변수, 시드 배정 규칙, 통계량 산출 방법을 기술한다.

\textbf{\chref{ch:results} 연구 결과}에서는 C0--C3 조건별 시뮬레이션 결과, 회귀분석,
수렴 진단, 가설 검증을 순차적으로 보고하고, 조건 간 효과 크기 분석,
회귀 모델의 설명력 구조, 이중채널 비선형 메커니즘, 전문가 검증,
선행연구와의 교차 정합성을 제시한다.

\textbf{\chref{ch:discussion} 논의 및 결론}에서는 연구 결과에 대한 종합 논의와
\chg{이론적, 방법론적, 실무적} 시사점,
6개 주요 산업 가이드라인(기준선 제외)을 제시하고,
연구 한계와 주요 성과를 정리하며 향후 연구 방향을 제시한다.%
}

\textbf{부록}은 세 개 장으로 구성하였다. 첫째는 기호·용어·수식 목록(부록 A), 둘째는 조직 성숙도 자가 진단 체크리스트(부록 B), 셋째는 전문가 검증 상세 설문 구성(부록 C)이다. 전체 시뮬레이션 코드와 대규모 재현 자료는 공개
재현 패키지에 수록하여 GitHub에 공개한다.

\tabref{tab:chapter_overview}는 \chg{이상에서 기술한 장과 부록의 구성, 그리고 각 장과 가설(H1--H10) 및 분석 조건(C0--C3)의 대응 관계를 요약한 것이다}.

\begin{table}[htbp]
  \centering
  \caption{논문 구성과 가설/조건 대응 관계}
  \label{tab:chapter_overview}
  \small
  \begin{tabularx}{\textwidth}{@{}c >{\raggedright\arraybackslash}p{2.4cm} X >{\centering\arraybackslash}p{1.5cm}@{}}
    \toprule
    장 & 제목 & 주요 내용 & 관련 가설/조건 \\
    \midrule
    1 & 서론 & 배경, 연구 필요성, 연구 목적, 범위 및 방법 & --- \\
    2 & 이론적 배경 & BCMS, 관측성, 불확실성, EVT, 바운딩 기법, 모델 설계, C0--C3, H1--\jrev{H10} & H1--\jrev{H10}, C0--C3 \\
    3 & 연구 방법 & 연구 설계, MC(몬테카를로 시뮬레이션) $n{=}10{,}000$ 실행 절차 & C0--C3 \\
    4 & 연구 결과 & 조건별 결과, 회귀, 수렴, 가설 검증, 교차분석, 전문가 검증 & H1--\jrev{H10}, C0--C3 \\
    5 & 논의 및 결론 & 종합 논의, 이론적/방법론적/실무적 시사점, 6개 주요 산업 가이드라인(기준선 제외), 한계, 성과 요약, 향후 연구 & 전체 \\
    \addlinespace
    부록 & 기호·수식·용어 일람, 전문가 평가 설문 문항, BCMS 성숙도 자가 진단 체크리스트 & & --- \\
    \bottomrule
  \end{tabularx}
\end{table}
